\section{Benedict XVI: Concessions and Doctrinal Dialogue}

Cardinal Ratzinger's role in the 1988 protocol gave his pontificate an unusually direct connection to the case. Benedict XVI received Bishop Bernard Fellay in August 2005. SSPX requests focused on broad freedom for the older Mass and removal of the 1988 excommunications as preliminary conditions for doctrinal discussion. These steps were not presented by the Holy See as a substitute for canonical reconciliation, but as obstacles that could be addressed to make it possible.

The 2007 motu proprio \emph{Summorum Pontificum} declared that the 1962 Roman Missal had never been juridically abrogated and established wider conditions for its use as the extraordinary form of the one Roman rite. The measure applied to the Latin Church broadly, not only to the SSPX. It nevertheless altered a central premise of the post-1988 landscape: attachment to the 1962 books could more readily be lived within ordinary canonical structures. The Society welcomed the recognition while continuing to argue that the crisis extended beyond liturgical permission.

\subsection{Remission is not recognition}

On 21 January 2009 the Congregation for Bishops remitted the excommunication of Bernard Fellay, Bernard Tissier de Mallerais, Richard Williamson, and Alfonso de Galarreta. The act responded to Fellay's petition and removed the declared censure from the four living bishops. It did not confer office, canonical mission, or status on the Society.

Benedict's letter to the bishops of 10 March made that boundary explicit. The remission concerned persons; the SSPX had no canonical status in the Church for doctrinal reasons, and its ministers did not legitimately exercise ministry. The pope also addressed the grave public controversy caused by Richard Williamson's denial of the scale and means of the Holocaust and acknowledged failures of information and communication. The episode joined a moral scandal to an already misunderstood juridical act and made the intended step toward dialogue appear to many observers as institutional rehabilitation.

\statusframe{Congregation for Bishops decree of 21 January and Benedict XVI's letter of 10 March 2009.}{The personal excommunications declared in 1988 were remitted for the four surviving bishops; doctrinal talks could proceed.}{The remission neither retroactively licensed the 1988 act nor granted the SSPX canonical status or legitimate ministry.}

\subsection{The doctrinal preamble}

\emph{Ecclesiae unitatem} of 2 July 2009 joined the \emph{Ecclesia Dei} Commission to the Congregation for the Doctrine of the Faith. The restructuring expressed the pope's judgment that the remaining problem was fundamentally doctrinal. Eight formal meetings followed between October 2009 and April 2011. Topics included the interpretation of the Council, the magisterium, liturgy, ecclesiology, religious liberty, and ecumenism.

In September 2011 the Congregation presented a doctrinal preamble whose acceptance would permit movement toward canonical recognition. Exchanges continued into 2012, and Roman authorities presented elements of a possible stable canonical solution. Public communiqués show alternating assessments of whether Fellay's responses were sufficient. No agreed doctrinal declaration and no canonical erection resulted.

The failure cannot be explained only as an administrative missed opportunity. The parties differed over the status of the contested conciliar teaching and the authority by which its interpretation would be settled. A formula flexible enough to permit continuing theological criticism still had to require a recognizable assent to the magisterium. The Society feared that signature would neutralize its foundational witness; Roman authorities feared that recognition without agreement would institutionalize rejection of the Council.

Internal dissent followed the negotiations. Williamson was excluded from the Society in 2012 after repeated disobedience. The dispute showed that even an unrealized agreement could destabilize a movement whose identity had been formed by refusing settlements judged insufficient.
